% ============================================================
% f2014_anhui_arts_q04_flowchart.tex
% 2014 安徽卷（文科）第 4 题：程序框图
% 来源 PDF：普通高考/2014/2014安徽文.pdf
% 原页：第 1 页（A3 横排）左栏，印面 859，第 4 题题干下方
% 结构：开始 → 赋初值 x=1、y=1 → z=x+y → 判断 z≤50；
%       是：x=y → y=z → 返回 z=x+y 之前；否：输出 z → 结束。
% ============================================================
\begin{tikzpicture}[>=Stealth, line width=0.7pt, font=\small,
  x=1cm, y=1cm, line cap=round, line join=round]

  % 开始
  \draw[rounded corners=6pt] (-0.95,0.42) rectangle (0.95,-0.42);
  \node at (0,0) {开始};
  \draw[->] (0,-0.42) -- (0,-0.85);

  % 赋初值
  \draw (-1.55,-0.85) rectangle (1.55,-1.70);
  \node at (0,-1.28) {$x=1$，$y=1$};
  \draw[->] (0,-1.70) -- (0,-2.10);

  % z=x+y
  \draw (-1.30,-2.10) rectangle (1.30,-2.95);
  \node at (0,-2.52) {$z=x+y$};
  \draw[->] (0,-2.95) -- (0,-3.40);

  % 判断框
  \draw (0,-3.40) -- (1.75,-4.25) -- (0,-5.10) -- (-1.75,-4.25) -- cycle;
  \node at (0,-4.25) {$z\leqslant 50$};

  % 否：向右输出
  \draw[->] (1.75,-4.25) -- (3.30,-4.25);
  \node at (2.20,-3.95) {否};
  \draw (3.30,-3.82) rectangle (5.50,-4.68);
  \node at (4.40,-4.25) {输出 $z$};
  \draw[->] (4.40,-4.68) -- (4.40,-5.05);

  % 结束
  \draw[rounded corners=6pt] (3.45,-5.05) rectangle (5.35,-5.90);
  \node at (4.40,-5.47) {结束};

  % 是：向左进入循环体
  \draw[->] (-1.75,-4.25) -- (-3.30,-4.25) -- (-3.30,-4.75);
  \node at (-2.30,-3.95) {是};
  \draw (-4.45,-4.75) rectangle (-2.15,-5.60);
  \node at (-3.30,-5.17) {$x=y$};

  \draw[->] (-3.30,-5.60) -- (-3.30,-5.95);
  \draw (-4.45,-5.95) rectangle (-2.15,-6.80);
  \node at (-3.30,-6.37) {$y=z$};

  % 回边：向下、向左、向上，再水平汇入 z=x+y 与判断框之间
  \draw (-3.30,-6.80) -- (-3.30,-7.25) -- (-6.30,-7.25) -- (-6.30,-3.18) -- (0,-3.18);
  \draw[->] (0,-3.18) -- (0,-3.40);

\end{tikzpicture}
